\documentclass[aspectratio=169,12pt]{beamer}

\usepackage[T1]{fontenc}
\usepackage[utf8]{inputenc}
\usepackage{lmodern}
\usepackage{amsmath,amssymb,booktabs}
\usepackage{appendixnumberbeamer}

% A plain mathematical Beamer layout. All formulas remain editable.
\definecolor{Ink}{HTML}{183447}
\definecolor{Teal}{HTML}{087D82}
\definecolor{Ochre}{HTML}{A25B12}
\definecolor{Muted}{HTML}{596974}
\setbeamercolor{normal text}{fg=Ink,bg=white}
\setbeamercolor{frametitle}{fg=Ink,bg=white}
\setbeamercolor{title}{fg=Ink}
\setbeamercolor{structure}{fg=Teal}
\setbeamercolor{alerted text}{fg=Teal}
\setbeamercolor{footline}{fg=Muted,bg=white}
\setbeamerfont{title}{size=\fontsize{28}{33}\selectfont,series=\bfseries}
\setbeamerfont{frametitle}{size=\fontsize{20}{24}\selectfont,series=\bfseries}
\setbeamerfont{subtitle}{size=\large}
\setbeamerfont{footline}{size=\scriptsize}
\setbeamersize{text margin left=11mm,text margin right=11mm}
\setbeamertemplate{navigation symbols}{}
\setbeamertemplate{itemize item}{\small\raisebox{0.15ex}{\(\bullet\)}}
\setbeamertemplate{frametitle}{\vspace{4mm}\insertframetitle\par\vspace{2mm}}
\setbeamertemplate{footline}{%
  \hspace*{11mm}\usebeamerfont{footline}Hodge theory\hfill
  \insertframenumber\hspace*{11mm}\vspace*{3mm}\par}
\setlength{\parskip}{0.5em}

\newcommand{\db}{\bar\partial}
\newcommand{\C}{\mathbb C}
\newcommand{\Z}{\mathbb Z}
\newcommand{\Lie}{\mathcal L}
\newcommand{\fouriermode}{e_{n,m}}
\newcommand{\zero}[1]{{\color{Ochre}#1}}
\newcommand{\aside}[1]{\par{\small\color{Muted}#1\par}}
\pdfstringdefDisableCommands{\def\translate#1{#1}}

\title{The \texorpdfstring{$\bar\partial$}{dbar}-Poincar\'e lemma}
\subtitle{One complex dimension}
\author{Daniel Gonz\'alez Casanova Azuela}
\date{}
\hypersetup{pdfauthor={Daniel Gonzalez Casanova Azuela},
  pdftitle={Lecture 12: The dbar-Poincare lemma in one complex dimension}}

\begin{document}

\begin{frame}[plain]
  \vspace{7mm}
  {\small\color{Teal}HODGE THEORY\par}
  \vspace{5mm}
  {\usebeamerfont{title}\inserttitle\par}
  \vspace{3mm}
  {\large\insertsubtitle\par}
  \vspace{6mm}
  {\normalsize\insertauthor\par}
  \vfill
  \aside{Based on Misha Verbitsky's lecture 12, IMPA, 2025.}
  \vspace{3mm}
\end{frame}

\begin{frame}{The equation we want to solve}
  Let $D\subset\C$ be a disk, with coordinate $z=x+iy$.

  Given a smooth $(0,1)$-form
  \[
    \eta=g(z,\bar z)\,d\bar z
  \]
  extending smoothly to a neighbourhood of $\overline D$, we want
  \[
    \alert{\db u=\eta},\qquad u\in C^\infty(D).
  \]
  In coordinates this is just
  \[
    \db u=\frac12(u_x+i u_y)\,d\bar z.
  \]
  \aside{We first solve the equation on a torus up to a constant term.}
\end{frame}

\begin{frame}{Functions on the square torus}
  We start with the elliptic curve
  \[
    E=\C/(2\pi\Z+2\pi i\Z).
  \]
  Think of $(0,2\pi]\times(0,2\pi]$ with opposite edges identified.

  A function on $E$ is a function on $\C$ satisfying
  \[
    f(x+2\pi,y)=f(x,y),\qquad
    f(x,y+2\pi)=f(x,y).
  \]
  \alert{There are two real coordinates, so there will be two frequencies.}
\end{frame}

% Sources: HP, functional-analysis.tex, example-fourier-orthonormal-system
% (use integer indices); Verbitsky, lecture 12, pp. 4--5.
\begin{frame}{The Fourier basis}
  The functions
  \[
    \fouriermode(x,y)=e^{i(nx+my)}=e^{inx}e^{imy},
    \qquad (n,m)\in\Z^2,
  \]
  form an orthonormal basis of $L^2(E)$ for the measure
  $dx\,dy/(2\pi)^2$.

  Therefore every $f\in L^2(E)$ has a basis expansion
  \[
    \alert{f=\sum_{n,m\in\Z}a_{n,m}\fouriermode},
    \qquad a_{n,m}=\langle f,\fouriermode\rangle.
  \]
  \aside{The equality means convergence in $L^2$. We use the Fourier basis theorem without proving it here.}
\end{frame}

\begin{frame}{Fourier modes of a form}
  Expand the coefficient of a smooth $(0,1)$-form:
  \[
    \eta=g\,d\bar z
      =\sum_{n,m\in\Z}\underbrace{a_{n,m}\fouriermode\,d\bar z}_{\eta_{n,m}}.
  \]
  The summand $\eta_{n,m}$ is its \alert{$(n,m)$-Fourier mode}.

  \[
    \eta=\zero{a_{0,0}\,d\bar z}
      +\sum_{(n,m)\ne(0,0)}\eta_{n,m}.
  \]
  \aside{The zero mode has constant coefficients. A nonzero frequency means $(n,m)\ne(0,0)$.}
\end{frame}

\begin{frame}{Differentiation preserves the frequency}
  Differentiate one basis function:
  \[
    d\fouriermode=\fouriermode(in\,dx+im\,dy),
  \]
  \[
    \db\fouriermode=\frac{in-m}{2}\,\fouriermode\,d\bar z.
  \]
  Every coefficient is a constant multiple of \alert{the same $\fouriermode$}.

  Thus $d$ and $\db$ preserve the pair $(n,m)$.
  We can solve the equation \alert{one frequency at a time}.
\end{frame}

\begin{frame}{Closed forms and exact forms}
  A form $\eta$ is \alert{$\db$-closed} if $\db\eta=0$.

  It is \alert{$\db$-exact} if $\eta=\db u$ for some $u$.

  In complex dimension one, every $(0,1)$-form is closed:
  \[
    \db(g\,d\bar z)
      =\frac{\partial g}{\partial\bar z}\,
        \underbrace{d\bar z\wedge d\bar z}_{0}=0.
  \]
  \aside{Our task is to find primitives. Closedness is already automatic in this degree.}
\end{frame}

% Source: Verbitsky, lecture 12, pp. 13--14.
% Conventions throughout this deck: z=x+iy, I(partial_x)=partial_y,
% dbar f=1/2(f_x+i f_y) dbar z. These fix the plus sign below.
% The formula is used here for constant translation fields on a flat torus.
\begin{frame}{The complex Cartan formula}
  Put $X=\partial_x$, so $IX=\partial_y$.

  For these constant vector fields,
  \[
    \alert{\db\iota_X+\iota_X\db
      =\frac12\bigl(\Lie_X+i\Lie_{IX}\bigr)}.
  \]
  Here $\iota_X$ is contraction. On a one-form, $\iota_X\eta=\eta(X)$.
  In particular,
  \[
    \iota_X(g\,d\bar z)=g.
  \]
  \aside{Here $\Lie_X$ and $\Lie_{IX}$ simply differentiate the coefficients in the $x$- and $y$-directions.}
\end{frame}

\begin{frame}{Primitives of nonzero modes}
  For $\eta_{n,m}=a_{n,m}\fouriermode\,d\bar z$, we have
  \[
    \Lie_X\eta_{n,m}=in\eta_{n,m},\qquad
    \Lie_{IX}\eta_{n,m}=im\eta_{n,m}.
  \]
  Since $\db\eta_{n,m}=0$, Cartan's formula gives
  \[
    \db(\iota_X\eta_{n,m})=\frac{in-m}{2}\,\eta_{n,m}.
  \]
  If $(n,m)\ne(0,0)$, then $in-m\ne0$, hence
  \[
    \alert{\eta_{n,m}
      =\db\left(\frac{2a_{n,m}}{in-m}\fouriermode\right)}.
  \]
\end{frame}

\begin{frame}{Adding the primitives}
  Define a function on the torus by
  \[
    f=\sum_{(n,m)\ne(0,0)}
       \frac{2a_{n,m}}{in-m}\fouriermode.
  \]
  Differentiating term by term gives
  \[
    \db f=\eta-\zero{a_{0,0}\,d\bar z}.
  \]
  Therefore
  \[
    \alert{\eta=\zero{c\,d\bar z}+\db f},\qquad c=a_{0,0}.
  \]
  \aside{Fourier fact used here: for smooth $g$, this series and all its derivatives converge uniformly, so $f$ is smooth.}
\end{frame}

\begin{frame}{The constant part in cohomology}
  Dolbeault cohomology identifies closed forms that differ by an exact form:
  \[
    H^{0,1}_{\db}(E)
      =\frac{\{\text{smooth closed $(0,1)$-forms}\}}
             {\{\db f:f\in C^\infty(E)\}}.
  \]
  Since $\eta=c\,d\bar z+\db f$,
  \[
    [\eta]=c[d\bar z],\qquad
    \alert{H^{0,1}_{\db}(E)\cong\C}.
  \]
  \aside{$\db f$ always has zero constant part: differentiation preserves frequencies and kills the constant term. Thus $c\,d\bar z$ is exact only when $c=0$.}
\end{frame}

\begin{frame}{The disk inside an elliptic curve}
  Choose $L>0$ so large that $\overline D$ lies in the interior
  of a fundamental square for
  \[
    E_L=\C/(L\Z+iL\Z).
  \]
  The quotient map embeds a neighbourhood of $\overline D$ into $E_L$.

  The same Fourier calculation works for any period $L$.
  \aside{Rescaling the coordinate by $w=2\pi z/L$ reduces it to the square torus we just used.}

  \alert{On this neighbourhood, $z$ and $\bar z$ are well-defined functions.}
\end{frame}

% Source: Verbitsky, lecture 12, p. 15. Multiplication by a cutoff is safe
% here because every (0,1)-form on a complex curve is dbar-closed.
\begin{frame}{Extending the form to the torus}
  Let $U$ be a neighbourhood of $\overline D$ inside that square,
  where $\eta$ is defined.

  Choose a smooth function $\chi$ that is $1$ near $\overline D$
  and has support contained in $U$.
  \[
    \widetilde\eta=
      \begin{cases}
        \chi\eta & \text{on }U,\\[1mm]
        0 & \text{outside }U.
      \end{cases}
  \]
  This is a smooth $(0,1)$-form on $E_L$, equal to $\eta$ on $D$.

  \aside{It is automatically $\db$-closed, just like every $(0,1)$-form on a complex curve.}
\end{frame}

\begin{frame}{The primitive on the disk}
  On the elliptic curve, our calculation gives
  \[
    \widetilde\eta=c\,d\bar z+\db f.
  \]
  On the disk, $\bar z$ is a function and $\db\bar z=d\bar z$.
  Restricting to $D$ therefore gives
  \[
    \alert{\eta=\db\bigl(c\bar z+f|_D\bigr)}.
  \]
  \vspace{2mm}
  This is the required primitive.

  \aside{On the whole torus, $\bar z$ is not periodic, so it cannot serve as a global primitive.}
\end{frame}

\appendix

\begin{frame}{Supplement: the $(1,1)$ case}
  A $(1,1)$-form on the disk is $\eta=g\,dz\wedge d\bar z$.

  By the $(0,1)$ case, choose $u$ with
  \[
    \db u=g\,d\bar z.
  \]
  Then
  \[
    \db(-u\,dz)=-\db u\wedge dz
      =g\,dz\wedge d\bar z=\eta.
  \]
  \aside{These are the two possible types with positive antiholomorphic degree in complex dimension one: $(0,1)$ and $(1,1)$.}
\end{frame}

\begin{frame}{Supplement: the constant representatives}
  The same calculation for each type gives
  \begin{center}
  \renewcommand{\arraystretch}{1.45}
  \begin{tabular}{ccc}
    \toprule
    Type & Constant representative & Cohomology on $E$ \\
    \midrule
    $(0,0)$ & $1$ & $H^{0,0}_{\db}(E)\cong\C$ \\
    $(1,0)$ & $dz$ & $H^{1,0}_{\db}(E)\cong\C$ \\
    $(0,1)$ & $d\bar z$ & $H^{0,1}_{\db}(E)\cong\C$ \\
    $(1,1)$ & $dz\wedge d\bar z$ & $H^{1,1}_{\db}(E)\cong\C$ \\
    \bottomrule
  \end{tabular}
  \end{center}
  \aside{Here $(p,q)$ records the number of $dz$ and $d\bar z$ factors. Each cohomology class has a unique constant representative.}
\end{frame}

\end{document}
